\myInitialCap{设}计模式把我们面向对象的软件设计技能提升到一个新水平。它们是模型，我们可以据此为许多常见的软件架构问题创建设计良好的解决方案。软件架构指的是我们如何组织应用的代码——我们所设计的类与子类，以及它们之间的相互关系和在运行时的交互方式。多年以来，各行各业的开发者已经证明，这些模式所刻画的是可靠而灵活、且遵循良好设计原则的解决方案。

以下是设计模式可以帮助解决的一些示例架构问题：

\begin{itemize}
\item 一个应用使用一族算法。我们能否封装每一个算法，并使它们可以互换？
\item 一个应用需要创建不同的对象集合。这些对象可以组织成若干类别或族。能否在对象创建上保持灵活，又不把来自不同族的对象混淆在一起？
\item 一个应用管理着若干不同的对象序列。应用能否逐个处理某个序列中的对象，而无需知道任何序列是如何实现的？
\item 应用中的一个对象在运行时产生数据流，供应用中的其他对象消费。我们能否这样设计应用：生产者无需知道消费者是谁、也无须知道每个消费者如何处理数据，并且能够在运行时添加或移除消费者？
\end{itemize}

这些章节所涵盖的设计模式，都建立在前面章节所介绍的面向对象概念和设计原则之上。这些模式的共同目标包括：封装变化之处，以及设计内聚且松耦合的类。设计模式可以减少硬编码的选项，使应用能够在运行时灵活地做出选择。

设计模式不是可以直接复制粘贴的源代码，也不是从库中导入的代码。我们通过把一个模式当作模型来改编它，编写旨在解决应用中某个特定架构问题的定制代码。这些模型涉及相互协作的类，与创建 vector、map、树或链表之类的基本数据结构无关。

设计模式不是被发明出来的，而是被发现出来的。多年来，经验丰富的开发者在许多软件项目中都遇到过常见的软件架构问题。从这些问题的最佳解决方案中，他们看出了类在设计、类之间的关系，以及对象在运行时交互上的模式。

经验丰富的程序员能够识别出可以用恰当的设计模式解决的架构问题，然后决定是否把该模式用作创建问题解决方案的模型。最早提出软件设计模式这一想法的四位计算机科学家写了一本书，描述了 23 个原始模式。本部分的各章将介绍其中的 13 个。

\begin{mySidebar}{设计模式与“四人组”}

设计模式这一概念最早是由克里斯托弗·亚历山大（Christopher Alexander）等人的著作 \emph{A Pattern Language: Towns, Building, Construction}（Oxford University Press，1977）推广开来的。这本书包含了用于解决实体建筑常见建筑问题的模式（门窗放在哪里、有多少层等），以及如何设计社区。

四位计算机科学家——埃里希·伽马（Erich Gamma）、理查德·赫尔姆（Richard Helm）、拉尔夫·约翰逊（Ralph Johnson）和约翰·弗利西斯（John Vlissides）——把这个概念加以改编，用于解决面向对象软件的常见架构问题。他们写了 \emph{Design Patterns: Elements of Reusable Object-Oriented Software}（Addison-Wesley，1995）。这本书常被称为\emph{“四人组”}（Gang of Four，GoF）之书。它包含 23 个软件设计模式，分为三大类：创建型模式、结构型模式和行为型模式。本章及随后几章将介绍其中 13 个最常见的模式。在边栏中，我直接从 GoF 书中引用了每个模式的意图，这些引用以（GoF \emph{nnn}）的记法标注出处，其中 \emph{nnn} 是该书的页码。

此后，众多作者发现并撰写了其他软件设计模式。这些新模式常用于专门的软件领域，例如网络与 Web、数据库和 GUI 编程。面向对象之外的编程范式，例如函数式编程，也有它们自己的设计模式。
\end{mySidebar}

\section*{设计模式的益处}

每当在开发应用时遇到软件架构问题，我们都不应假定这个问题是独一无二的。如果能够使用某个设计模式，我们就不必重新发明轮子。

如果所有应用开发者都精通设计模式，设计会议就会高效得多。例如，在会议中，一位开发者可能会说：“我们应该用策略设计模式（Strategy Design Pattern）来解决这个问题。”会议中的其他每位开发者都会明白这句话的含义，以及代码应当如何编写。设计模式是一种设计词汇表，它让会议保持在较高的层次上，而不会被如何实现解决方案的编码细节拖住。一个模式规定了解决方案中各类之间的关系，以及它们的对象在运行时的交互方式。一个设计良好的应用可以由多个设计模式来建模。

\section*{设计模式的讲解方式}

以下各章按照如下大纲来讲解每个设计模式：

\begin{enumerate}
\item 使用设计模式之前 a. 描述该设计模式所解决的软件架构问题 b. 描述一个存在该架构问题的简短应用 c. 展示该问题的 UML 图和“之前”的应用代码 d. 列出在使用该设计模式之前，那个解决方案存在的缺陷
\item 使用设计模式之后 a. 依据该设计模式建模得到的 UML 图和“之后”的应用代码 b. 列出第二个解决方案的益处
\item 设计模式的通用模型 a. 该设计模式通用模型的 UML 图 b. 一张表，把该模式通用模型的各个组成部分，与我们对该应用的第二个解决方案的各个组成部分对应起来
\end{enumerate}

如果一章讲解多个设计模式，并且示例应用是前一个模式中步骤 e 的解决方案的延伸，则省略步骤 1.c 和 1.d。

一个常见的问题是：“我怎么知道何时该使用某个设计模式？”要认识到何时使用模式、以及该用哪一个，需要一些经验与实践。

设计模式不是一个具体的东西，而是一个抽象模型，我们可以据此为某个软件架构问题创建定制解决方案。因此，为了讲解这些模式，本部分的各章包含许多“之前”与“之后”的 UML 图和代码示例。\emph{仔细研究这些示例，理解“之前”解决方案的缺陷，以及依据每个模式建模得到的“之后”解决方案的益处，这一点很重要。}

\section*{体育示例}

设计模式的应用示例涉及一所学校的体育部、棒球、排球和橄榄球这三项运动，以及生成有关比赛的报告。对于不熟悉棒球和排球的读者，下面高度简化地说明这两项比赛的打法与计分方式。

\begin{mySidebar}{棒球}

棒球由两支球队进行，每队九名球员。两队轮流进攻（击球）和防守（守备）。比赛在一个包含棒球内场的球场上进行，内场一个角是本垒，逆时针方向的其他三个角依次是一垒、二垒和三垒。内场各边长 90 英尺（约 27 米）。球场延伸到内场之外。

进攻方按事先排定的顺序，一次派出一名球员担任击球员，手持球棒站在本垒旁。防守方指定一名球员为投手，站在内场中央，把棒球投过本垒，击球员则试图用球棒击中它。

成功击中球的击球员会试图逆时针绕垒奔跑。视对方守备的好坏以及击球员跑垒的快慢，击球员可以安全到达一垒（击出一垒安打）、二垒（二垒安打）、三垒（三垒安打），或一路跑完（击出全垒打）。当前在各垒上的球员也会试图前进。每有一名球员跑回本垒得分，该队就得到一分（一个得分）。

如果在数次尝试之后，击球员未能击中投手投出的球，则出局。击中了球但未能成功到达某个垒的击球员同样出局。进攻方球员每累计三次出局，两队就交换角色。每队各轮一次进攻和防守，构成一局，打满九局后比赛结束。因此，在这个高度简化的比赛版本中，当每队各自出局 27 次后，比赛就结束了。
\end{mySidebar}

\begin{mySidebar}{排球}

排球由两支球队进行，每队六名球员，球场为长方形，中间隔着一张约 2.5 米高的球网。

一回合的比赛从某队的一名球员用手把球击过球网发球开始。此后，两队把球在网的两侧来回击打，直到有一队未能接住，从而送给对方一分并结束该回合。

当某队率先得到 15 分时，比赛结束。
\end{mySidebar}
